\hypertarget{appendices}{%
\chapter*{Appendices}\label{appendices}}
\addcontentsline{toc}{chapter}{Appendices}

\hypertarget{appendix-a-synthetic-data-generation-full-parameters}{%
\section*{Appendix A: Synthetic Data Generation, Full
Parameters}\label{appendix-a-synthetic-data-generation-full-parameters}}
\addcontentsline{toc}{section}{Appendix A: Synthetic Data Generation,
Full Parameters}

This appendix gives the generator's exact parameters, complementing the
methodology and justification given in Chapter 3. The full generator
source is \texttt{data/generate\_synthetic\_bids.py}.

\textbf{Table A.1: Feature marginal distributions}

\begin{longtable}[]{@{}
  >{\raggedright\arraybackslash}p{(\columnwidth - 2\tabcolsep) * \real{0.5000}}
  >{\raggedright\arraybackslash}p{(\columnwidth - 2\tabcolsep) * \real{0.5000}}@{}}
\toprule\noalign{}
\begin{minipage}[b]{\linewidth}\raggedright
Feature
\end{minipage} & \begin{minipage}[b]{\linewidth}\raggedright
Values and probabilities
\end{minipage} \\
\midrule\noalign{}
\endhead
\bottomrule\noalign{}
\endlastfoot
Client segment & AAA publisher 30\%, mid-size publisher 32\%, indie
studio 28\%, platform holder 10\% \\
Project type & Full development 18\%, co-development 34\%, porting 20\%,
live operations 16\%, art services 12\% \\
Relationship history & New client 45\%, prior/poor delivery 10\%,
prior/neutral delivery 22\%, prior/good delivery 23\% \\
Incumbency position & Greenfield 55\%, defending existing work 22\%,
displacing a competitor 23\% \\
\end{longtable}

Deal origin, pre-RFP engagement, effort, price position, requirement
coverage, portfolio similarity, geo/timezone fit, security credentials,
and competing-bidder count are each drawn conditionally or from
continuous distributions (lognormal for effort; Beta for geo/timezone
fit; Poisson for competing bidders) as described in Chapter 3; exact
distributional parameters are given in the generator source.

\textbf{Table A.2: Encoded log-odds weights (relative effect on win
propensity)}

\begin{longtable}[]{@{}
  >{\raggedright\arraybackslash}p{(\columnwidth - 2\tabcolsep) * \real{0.5000}}
  >{\raggedright\arraybackslash}p{(\columnwidth - 2\tabcolsep) * \real{0.5000}}@{}}
\toprule\noalign{}
\begin{minipage}[b]{\linewidth}\raggedright
Feature / level
\end{minipage} & \begin{minipage}[b]{\linewidth}\raggedright
Encoded effect
\end{minipage} \\
\midrule\noalign{}
\endhead
\bottomrule\noalign{}
\endlastfoot
Relationship history: new / poor prior / neutral prior / good prior &
−0.35 / −0.55 / +0.35 / +0.95 \\
Deal origin: inbound & +0.75 \\
Pre-RFP engagement & +0.65 \\
Incumbency: greenfield / defending / displacing & 0 / +0.55 / −0.60 \\
Price position & −0.045 per point above benchmark; −0.017 per point
below \\
Requirement coverage & +2.4 × (coverage − 0.75) \\
Portfolio similarity & +1.3 × (similarity − 0.5) \\
Effort & −0.0012 × (person-months − 90) \\
Geo/timezone fit & +0.9 × (fit − 0.6) \\
Security credentials & +0.20 (base; see interaction below) \\
Competing bidders & −0.12 × (count − 3) \\
Interaction: weak portfolio (\textless{} 0.45) × new client & −1.1 \\
Interaction: security credentials × AAA/platform segment & +0.85 \\
Interaction: price overage × indie-studio segment & −0.03 per point \\
Latent noise & Gaussian, σ = 1.25 \\
Label noise & 5\% of outcomes flipped post-draw \\
\end{longtable}

The intercept is calibrated by binary search so the realised win rate
lands close to 30\%; on the run this thesis reports, the realised win
rate was 30.6\% across 500 bids.

\hypertarget{appendix-b-power-bi-migration-plan}{%
\section*{Appendix B: Power BI Migration
Plan}\label{appendix-b-power-bi-migration-plan}}
\addcontentsline{toc}{section}{Appendix B: Power BI Migration Plan}

This appendix sets out concretely what porting the delivered pipeline's
output to Power BI, the originally proposed delivery mechanism (Section
1.6, Section 3.4), would involve, since the pipeline was deliberately
built to make this a presentation-layer addition rather than a rebuild.

\begin{enumerate}
\def\labelenumi{\arabic{enumi}.}
\tightlist
\item
  \textbf{Data import.} \texttt{outputs/model\_results.json}, or the
  nine flattened per-section tables already generated for the Claude
  Artifact dashboard (\texttt{dashboard/tables/*.json}, Chapter 4),
  would be imported via Power Query as Power BI tables; the flattened
  tables are directly usable as-is, since Power BI, like the Claude
  Artifact platform, works natively with tabular data.
\item
  \textbf{Report pages.} Each dashboard section described in Chapter 4
  maps to one Power BI report page: a performance-summary page (card
  visuals for AUC, Brier score, and the Spearman correlation, plus a
  line chart for the ROC curve and a matrix visual for the confusion
  matrix); a SHAP importance page (a horizontal bar chart); a
  hypothesis-comparison page (a table visual with conditional formatting
  on the rank-shift column); and an example-bids page (a table with
  drill-through to a detail page per bid).
\item
  \textbf{The scorer.} This is the one component that cannot be ported
  directly, since Power BI's native interactivity does not support a
  live, client-side arithmetic scorer in the way the web artefact's
  JavaScript implementation does (Section 4.4). The closest native
  equivalent is a set of Power BI ``What if'' parameters, one slider per
  numeric feature and one dropdown per categorical feature, feeding a
  DAX measure that replicates the logistic surrogate's exact computation
  (intercept plus the sum of each standardised feature times its
  coefficient, passed through a logistic transform), using the
  coefficients already exported in
  \texttt{outputs/model\_results.json}'s \texttt{scorer\_model} object.
  This is a direct, mechanical translation of existing, already-computed
  numbers into DAX, not new modelling work.
\item
  \textbf{Refresh.} A scheduled Power BI refresh, pulling an updated
  \texttt{model\_results.json} each time the underlying pipeline is
  re-run (for example, on a periodic schedule once real bid data
  accumulates, Chapter 7), would restore the MLOps-style automated
  refresh specified in the original proposal, implemented as a scheduled
  task that re-runs \texttt{src/train\_model.py} and republishes its
  output file to wherever Power BI's refresh step reads from.
\end{enumerate}

\hypertarget{appendix-c-repository-and-pipeline-structure}{%
\section*{Appendix C: Repository and Pipeline
Structure}\label{appendix-c-repository-and-pipeline-structure}}
\addcontentsline{toc}{section}{Appendix C: Repository and Pipeline
Structure}

\textbf{Table C.1: Pipeline scripts, their CRISP-DM stage (Chapter 3),
and how to run them}

{\footnotesize
\begin{longtable}[]{@{}
  >{\raggedright\arraybackslash}p{(\columnwidth - 6\tabcolsep) * \real{0.3000}}
  >{\raggedright\arraybackslash}p{(\columnwidth - 6\tabcolsep) * \real{0.1700}}
  >{\raggedright\arraybackslash}p{(\columnwidth - 6\tabcolsep) * \real{0.2300}}
  >{\raggedright\arraybackslash}p{(\columnwidth - 6\tabcolsep) * \real{0.3000}}@{}}
\toprule\noalign{}
\begin{minipage}[b]{\linewidth}\raggedright
Script
\end{minipage} & \begin{minipage}[b]{\linewidth}\raggedright
CRISP-DM stage
\end{minipage} & \begin{minipage}[b]{\linewidth}\raggedright
Output
\end{minipage} & \begin{minipage}[b]{\linewidth}\raggedright
Run
\end{minipage} \\
\midrule\noalign{}
\endhead
\bottomrule\noalign{}
\endlastfoot
\texttt{data/generate\_synthetic\_bids.py} & Data understanding, data
preparation & \texttt{data/synthetic\_bids.csv} &
\texttt{python3\ data/generate\_synthetic\_bids.py} \\
\texttt{src/train\_model.py} & Modelling, evaluation &
\texttt{outputs/model\_results.json} &
\texttt{python3\ src/train\_model.py} \\
\texttt{dashboard/make\_tables.py} & Deployment &
\texttt{dashboard/tables/*.json}, \texttt{docs/data/*.json} &
\texttt{python3\ dashboard/make\_tables.py} \\
\texttt{dashboard/index.html} & Deployment & Claude Artifact dashboard &
Published as a hosted Artifact \\
\texttt{docs/index.html} & Deployment & Static GitHub Pages site &
Served from the repository's \texttt{docs/} folder \\
\end{longtable}
}

\texttt{data/methodology\_synthetic\_data.md} gives the synthetic-data
methodology in its original, standalone form, written before this thesis
and incorporated into Chapter 3. All code is version-controlled in the
project's Git repository alongside this thesis document, consistent with
Hevner et al.'s \citeyearpar{hevner2004design} requirement that the
research process be rigorously and repeatably documented (Chapter 3).

\hypertarget{appendix-d-feature-glossary}{%
\section*{Appendix D: Feature
Glossary}\label{appendix-d-feature-glossary}}
\addcontentsline{toc}{section}{Appendix D: Feature Glossary}

This glossary reproduces, for reference, the plain-language feature
definitions built into the artefact's clickable glossary (Chapter 4),
covering every feature used in modelling.

\begin{longtable}[]{@{}
  >{\raggedright\arraybackslash}p{(\columnwidth - 2\tabcolsep) * \real{0.5000}}
  >{\raggedright\arraybackslash}p{(\columnwidth - 2\tabcolsep) * \real{0.5000}}@{}}
\toprule\noalign{}
\begin{minipage}[b]{\linewidth}\raggedright
Feature
\end{minipage} & \begin{minipage}[b]{\linewidth}\raggedright
Definition
\end{minipage} \\
\midrule\noalign{}
\endhead
\bottomrule\noalign{}
\endlastfoot
History with this client & Whether this client is new to Virtuos, or a
repeat client, and if a repeat client, how the last project went:
poorly, fine, or well. \\
Who started the conversation & Whether the client approached Virtuos
about this opportunity (inbound), or Virtuos pursued the client
(outbound). \\
Talked before the formal bid & Whether Virtuos had conversations with
the client about the project before the official request for proposal
went out. \\
Starting position on this deal & Whether this is a brand-new opportunity
with no prior vendor, work Virtuos is already doing that it is
defending, or work currently done by a competitor that Virtuos is trying
to win away. \\
Project size & The estimated effort for the project, in
person-months. \\
Price vs.~the market rate & How Virtuos's quoted price compares to the
typical rate for similar work; positive means priced above that
benchmark. \\
How much of the brief we covered & The share of the client's stated
requirements that the proposal fully addresses. \\
Relevant past work & How closely Virtuos's existing portfolio matches
the kind of work this project needs. \\
Time zone overlap & How much the working hours of Virtuos's team overlap
with the client's. \\
Security credentials in place & Whether Virtuos already holds the data
and IP security certifications this client's industry typically
requires. \\
Number of competing bidders & How many other studios were bidding for
the same project (flagged in Chapter 3 as a possible real-world leakage
risk, since it is often only known after the decision is made). \\
Type of client & AAA publisher, mid-size publisher, indie studio, or
platform holder. \\
Type of project & Full development, co-development, porting, live
operations, or art services. \\
\end{longtable}
